\newpage
\section*{ABSTRACT}  
\phantomsection
\addcontentsline{toc}{section}{ABSTRACT}

The precise positioning of the directional antenna is important for establishing stable wireless connections, especially in inexpensive and constrained systems where expensive phased arrays and beamforming technology cannot be implemented. This work presents the design and implementation of an autonomous antenna alignment system where the direction and orientation of an antenna can be controlled by a policy learned via reinforcement learning based on the RSSI values. The system consists of one stationary transmitter node and one receiver node consisting of an ESP32-S3, 2.4~GHz  high-gain directional antenna, azimuth stepper and tilt servo actuators. The measured RSSI value is averaged and encoded into a discretized state space to be used to retrieve pre-trained tabulated Q-learning policy on the ESP32-S3. Selected actions are then used to drive the azimuth and tilt actuators and control antenna direction accordingly. In order to evaluate the viability of reinforcement learning in such applications, the method is intended to be evaluated in relation to several baseline strategies, such as fixed direction, full angular scanning and hill climbing based on RSSI. The evaluation methodology will take into consideration the convergence speed, accuracy of the alignment, stability of the connection and mechanical efficiency of the antenna actuation both indoor and outdoor. When tested in a number of realistic environments, the reinforcement learning-based antenna alignment provided the average final RSSI of $-48.95$~dBm, while the exhaustive angular scanning achieved an RSSI of $-49.55$~dBm, implying only slightly lower maximum RSSI with respect to the exhaustive scanning. Nevertheless, the reinforcement learning reduced average convergence time from 54.20~s to 9.62~s, improving the speed of finding the best antenna orientation by almost 82.25\%.

\vspace{0.5cm}

\textit{Keywords: Directional Antenna Alignment, Embedded Systems, Policy-based Control, Reinforcement Learning, RSSI-based Optimization, Wireless Communication}
